Table 1 gives ecliptic longitude ($\lambda$) and parallax ($p$) at different
times ($t$), for a certain hypothetical minor planet. The baseline for the
parallax is the diameter of the earth. The time is expressed in years and for
your reference ecliptic longitudes of the Sun ($\lambda_\odot$) for the same
dates are also given the table. Let us assume that the orbital inclination of
this minor planet, with respect to the ecliptic, is negligible and the
eccentricity of the Earth's orbit is negligible.

\begin{description}
\item[(a)] Calculate the coordinates of the minor planet in the heliocentric
  polar coordinate system and put them in a polar plot. The x-axis in the plot
  should be directed towards the initial position of the minor planet. Draw
  the major axis of the orbit of the minor planet.\\
  Identify erroneous observation(s), if any. \hfill[38 points]
\item[(b)] Assuming the heliocentric orbit of the minor planet to be
  elliptical, determine
  \begin{description}
  \item[(i)] the semimajor axis length $a_p$.
  \item[(ii)] eccentricity $e$.
  \item[(iii)] the period $P$.
  \end{description}
  \hfill[6 points]
\item[(c)] Estimate the errors in the values of $P$, $a_p$ and $e$.
  \hfill[6 points]
\end{description}

\begin{center}
Table 1: Minor planet data

\begin{tabular}{cccc}
\hline
$t$ & $\lambda$ & $\lambda_\odot$ & $p$ \\
{[year]} & {[$^\circ$]} & {[$^\circ$]} & {[$''$]} \\
\hline
2012.3 & 336.73 & 40.95 & 3.82 \\
2012.6 & 3.44 & 134.83 & 7.24 \\
2012.9 & 50.71 & 242.08 & 7.09 \\
2013.4 & 94.52 & 64.84 & 2.40 \\
2013.6 & 121.40 & 134.59 & 2.16 \\
2013.9 & 154.31 & 241.82 & 2.75 \\
2014.2 & 25.33 & 353.29 & 3.16 \\
2014.5 & 148.51 & 99.04 & 1.99 \\
2014.8 & 176.26 & 205.45 & 1.83 \\
2015.0 & 216.33 & 280.19 & 2.03 \\
2015.3 & 187.5 & 28.55 & 2.897 \\
\hline
\end{tabular}
\end{center}
